\documentclass[11pt,a4paper]{article}
\usepackage[T1]{fontenc}
\usepackage{lmodern,amsmath,amssymb,amsthm,mathtools,booktabs,microtype}
\usepackage[margin=27mm]{geometry}
\usepackage[colorlinks=true,linkcolor=blue,citecolor=blue,urlcolor=blue]{hyperref}
\newtheorem{theorem}{Theorem}[section]
\newtheorem{lemma}[theorem]{Lemma}
\newtheorem{proposition}[theorem]{Proposition}
\newtheorem{corollary}[theorem]{Corollary}
\theoremstyle{remark}\newtheorem{remark}[theorem]{Remark}
\newcommand{\Q}{\mathbb Q}\newcommand{\Z}{\mathbb Z}
\newcommand{\F}{\mathbb F}
\DeclareMathOperator{\Tr}{Tr}
\setlength{\emergencystretch}{2em}
\title{Rational Combinations of Two Contact Values}
\author{}\date{10 November 2026}
\begin{document}\maketitle
\begin{abstract}
A single rational combination of two independent contact values
cannot have odd valuation whenever either contact value does.
We prove this statement on the exact eighth-power scale used in
the thirteen-variable construction under weak \(abc\).
The counterexamples have positive slopes that are sums of two
integer squares, nonintegral rational inputs, and an input residual
that is a unit at the prime in question.
The obstruction concerns this uniform local parity implication;
it gives no lower bound for arbitrary Diophantine definitions.
\end{abstract}

\section{The contact function and the statement}

For a prime \(p\), write \(v_p\) for the valuation on
\(\Q_p^\times\) normalized by \(v_p(p)=1\). Put
\begin{equation}\label{eq:R}
 R(t)=
 \frac{t(t+1)(2t+1)(3t+1)(4t+1)(5t+2)}
      {120(t-1)^8}.
\end{equation}
Its six simple zeros are
\[
 0,\quad -1,\quad-\frac12,\quad-\frac13,\quad-\frac14,
 \quad-\frac25.
\]
Fix an odd integer \(B\ge3\), an integer \(n\ne0\), and a rational
number \(T>6\). For \(X\in\Q\), define
\begin{equation}\label{eq:scale}
       \beta(X)=2X^2+B,\qquad a(X)=\beta(X)^8+T^8.
\end{equation}
In the application in~\cite{Conditional13}, one takes
\(T=T(n)=\Tr_{\Q(\sqrt2)/\Q}x([n]P)\).
The argument below uses only that this is a fixed rational number
greater than \(6\); it uses no \(abc\) inequality.

\begin{theorem}\label{thm:main}
Let \(\Phi\in\Q(Z_0,Z_1)\) be nonzero, and let \(S\) be any finite
set of primes. There are infinitely many primes \(p\notin S\)
with the following property.

There are positive integers \(\mu,\nu\), at least one of them odd,
and rational numbers \(X,\lambda\) such that
\begin{gather}
 X\in\tfrac12\Z\setminus\Z,\qquad
 \lambda=s^2+t^2>1\quad(s,t\in\Z),                                      \label{eq:domain}\\
 v_p(R(\lambda))=\mu,\qquad v_p(R(a(X)\lambda))=\nu,                    \label{eq:depths}\\
 \Phi(R(\lambda),R(a(X)\lambda))\ne0
 \quad\hbox{is defined},                                               \label{eq:defined}\\
 v_p\!\left(\Phi(R(\lambda),R(a(X)\lambda))\right)\in2\Z,\qquad
 v_p(\beta(X)-n)=0.                                                     \label{eq:miss}
\end{gather}
For each such \(p\), both \(\mu\) and \(\nu\) can be made arbitrarily
large. After fixing the depths, \(\lambda\) and \(|X|\) can also be
made arbitrarily large.
\end{theorem}

Thus the implication
\[
 \bigl(v_p(R(\lambda))>0\text{ odd}\ \text{or}\
       v_p(R(a(X)\lambda))>0\text{ odd}\bigr)
 \ \Longrightarrow\
 v_p\bigl(\Phi(R(\lambda),R(a(X)\lambda))\bigr)\text{ odd}
\]
cannot hold uniformly on this parameter domain, even after
discarding any prescribed finite set of primes.
Since \(X=q/2\) with \(q\) odd in the theorem,
\(\beta(X)=q^2/2+B\) is itself nonintegral.

\section{A valuation lemma}

The elementary obstruction is slightly more general than the
case of two contact values.

\begin{lemma}\label{lem:parity}
Let \(1\le m<d\), and let
\(\Phi_1,\ldots,\Phi_m\in\Q(Z_1,\ldots,Z_d)\) be nonzero rational
functions. There is a finite set of primes \(S_0\) and a sequence
\[
        w_N=(w_{N,1},\ldots,w_{N,d})\in\Z_{>0}^d
\]
such that every coordinate tends to infinity, at least one
coordinate of every \(w_N\) is odd, and the following holds.
For every \(p\notin S_0\) and every
\(z_1,\ldots,z_d\in\Q_p^\times\) with
\(v_p(z_i)=w_{N,i}\), all the \(\Phi_j(z_1,\ldots,z_d)\)
are defined and nonzero and have even valuation.
\end{lemma}

\begin{proof}
Write \(\Phi_j=P_j/Q_j\) with nonzero polynomials
\(P_j,Q_j\in\Z[Z_1,\ldots,Z_d]\).
Exclude the finitely many primes dividing a nonzero coefficient
of any of these polynomials.

For each polynomial, compare the linear forms
\(\langle e,w\rangle\) belonging to its distinct exponent vectors
\(e\). The equalities between these forms lie on finitely many
proper hyperplanes. Choose \(w\) in the positive orthant off all
these hyperplanes. There is an open cone \(C\) containing \(w\)
on which each \(P_j\) and each \(Q_j\) has a unique monomial of
smallest weight. Let their exponent vectors be \(u_j\) and
\(v_j\), respectively. If the valuation vector of \(z\) lies in
\(C\), the unique term of smallest valuation cannot cancel.
Consequently
\begin{equation}\label{eq:linear}
         v_p(\Phi_j(z))=\langle u_j-v_j,v_p(z)\rangle.
\end{equation}
This also proves that both evaluated polynomials are nonzero.

The matrix with rows \(u_j-v_j\), reduced modulo \(2\), has
\(m<d\) rows. Its kernel over \(\F_2\) contains a nonzero vector
\(\epsilon\in\{0,1\}^d\).
Choose a positive integer vector \(v\in C\), and set
\[
                         w_N=2Nv+\epsilon .
\]
For all sufficiently large \(N\), this vector lies in \(C\).
Every coordinate tends to infinity, its parity vector is the
fixed nonzero vector \(\epsilon\), and all the expressions
in~\eqref{eq:linear} are even. Discarding the finitely many initial
values of \(N\) proves the statement.
\end{proof}

The lemma concerns valuations of rational functions of \(d\)
independent values. Its linear-algebra statement is not a
variable-count lower bound for arbitrary systems of polynomial
equations.

\section{Simple roots of the prescribed scale}

We give the elementary prime-divisor argument needed to retain
the exact scale~\eqref{eq:scale}.

\begin{lemma}\label{lem:primevalues}
Let \(G\in\Z[X]\) be nonconstant with \(G(0)\ne0\).
There are infinitely many primes dividing nonzero values \(G(k)\),
\(k\in\Z\).
If \(G\) is separable over \(\Q\), then outside any finite set
there are infinitely many primes at which \(G\) has a simple
root modulo \(p\).
\end{lemma}

\begin{proof}
Put \(c=G(0)\). Given a finite set of primes, let \(M\) be their
product, with \(M=1\) for the empty set. The polynomial
\[
                        H(Y)=G(cMY)/c
\]
has integer coefficients, is nonconstant, has constant term
\(1\), and is congruent to \(1\) modulo every prime dividing
\(M\). Choose an integer \(k\) with \(|H(k)|>1\).
A prime divisor \(p\) of \(H(k)\) lies outside the specified set
and divides the nonzero integer \(G(cMk)\).
This proves infinitude.

If \(G\) is separable, enlarge the excluded set by the primes
dividing its leading coefficient and its nonzero discriminant.
A root of \(G\) modulo any remaining prime is simple.
\end{proof}

\begin{lemma}\label{lem:scale}
There are infinitely many primes \(p\), outside any prescribed
finite set, and \(\xi\in\Z_p\) such that
\[
       a(\xi)=1,\qquad a'(\xi)\in\Z_p^\times,\qquad
       \beta(\xi)-n\in\Z_p^\times .
\]
\end{lemma}

\begin{proof}
The polynomial
\[
               f(X)=(2X^2+B)^8+T^8-1
\]
has derivative \(32X(2X^2+B)^7\).
A common zero of \(f\) and its derivative would have either
\(X=0\) or \(2X^2+B=0\). These alternatives give, respectively,
\[
                    B^8+T^8-1=0,\qquad T^8-1=0,
\]
both impossible because \(B\ge3\) and \(T>6\).
Thus \(f\) is separable and \(f(0)\ne0\).
After clearing denominators, apply Lemma~\ref{lem:primevalues},
also excluding primes dividing the clearing factor and the
denominator of \(T\). Hensel's lemma lifts every resulting
simple root to a root \(\xi\in\Z_p\) with \(a'(\xi)\) a unit.

The nonzero rational number \(n^8+T^8-1\) is positive.
Exclude the primes dividing its numerator. If
\(\beta(\xi)\equiv n\pmod p\), the equation
\(a(\xi)=1\) would give \(n^8+T^8-1\equiv0\pmod p\), a
contradiction. The enlarged excluded set is still finite.
\end{proof}

\section{Realizing the two contact valuations}

At the common contact \(-1\), direct differentiation gives
\begin{equation}\label{eq:derivative}
                            R'(-1)=-\frac{3}{5120}.
\end{equation}
Hence, for \(p\notin\{2,3,5\}\) and \(z\equiv-1\pmod p\),
\begin{equation}\label{eq:simple}
                         v_p(R(z))=v_p(z+1).
\end{equation}

\begin{lemma}\label{lem:norm}
Let \(p\) be odd and let \(\lambda_*\in\Z_p\) satisfy
\(\lambda_*\equiv-1\pmod p\). For every positive integer \(M\)
there exist integers \(s,t\) such that
\[
                  s^2+t^2\equiv\lambda_*\pmod{p^M}.
\]
The integer \(s^2+t^2\) can be chosen arbitrarily large.
\end{lemma}

\begin{proof}
The sets of residues \(u^2\) and \(-1-v^2\) modulo \(p\)
both have cardinality \((p+1)/2\), so they intersect.
Thus \(u^2+v^2=-1\) has a solution over \(\F_p\), and its
two coordinates cannot both be zero. Interchange them if
necessary so that \(u\ne0\). Fix an integer lift \(t\) of \(v\).
Hensel's lemma, applied to
\(S^2+t^2-\lambda_*\), gives a root \(s_*\in\Z_p^\times\).
Any integer \(s\equiv s_*\pmod{p^M}\) gives the required
congruence. Adding arbitrary integer multiples of \(p^M\)
to \(s\) makes the norm arbitrarily large.
\end{proof}

\begin{lemma}\label{lem:independent}
For each prime furnished by Lemma~\ref{lem:scale} outside
\(\{2,3,5\}\), and for every pair of positive integers
\(\mu,\nu\), there are \(X,\lambda\) satisfying
\eqref{eq:domain}, \eqref{eq:depths}, and
\(v_p(\beta(X)-n)=0\).
The quantities \(\lambda\) and \(|X|\) can be chosen arbitrarily
large after fixing \(\mu,\nu\).
\end{lemma}

\begin{proof}
Choose \(\xi\) as in Lemma~\ref{lem:scale}, and put
\[
 \lambda_*=-1+p^\mu,\qquad
 \tau_*=-1+p^\nu,\qquad
 a_*=\tau_*/\lambda_*.
\]
Both \(\lambda_*\) and \(\tau_*\) are units, and \(a_*\equiv1\pmod p\).
Since \(a'(\xi)\) is a unit, Hensel's lemma gives
\(X_*\equiv\xi\pmod p\) with \(a(X_*)=a_*\).
Equivalently, the two-variable map
\[
            (X,\lambda)\longmapsto
             \bigl(R(\lambda),R(a(X)\lambda)\bigr)
\]
has, at \((\xi,-1)\), Jacobian matrix
\[
 \begin{pmatrix}
 0 & R'(-1)\\
 -a'(\xi)R'(-1) & R'(-1)
 \end{pmatrix},
 \qquad
 \det=R'(-1)^2a'(\xi)\in\Z_p^\times .
\]

Choose \(M>\max(\mu,\nu)\).
By Lemma~\ref{lem:norm}, take
\(\lambda=s^2+t^2>1\) congruent to \(\lambda_*\) modulo \(p^M\).
The Chinese remainder theorem gives an odd integer \(q\) such that
\[
                       q\equiv2X_*\pmod{p^M}.
\]
Set \(X=q/2\). Then \(X\in\tfrac12\Z\setminus\Z\) and
\(X\equiv X_*\pmod{p^M}\).
The coefficients of \(a\) are \(p\)-integral, so
\[
       \lambda\equiv-1+p^\mu\pmod{p^M},\qquad
       a(X)\lambda\equiv-1+p^\nu\pmod{p^M}.
\]
Their respective contact depths are exactly \(\mu\) and \(\nu\);
\eqref{eq:simple} proves~\eqref{eq:depths}.
Moreover \(X\equiv\xi\pmod p\), so \(\beta(X)-n\) remains a unit.
Arbitrarily large lifts of \(q\) modulo \(2p^M\), independently
of the arbitrarily large norms in Lemma~\ref{lem:norm}, give
the last assertion.
\end{proof}

\begin{proof}[Proof of Theorem~\ref{thm:main}]
Apply Lemma~\ref{lem:parity} with \(d=2\), \(m=1\), and the
given \(\Phi\). Add its finite exceptional set and
\(\{2,3,5\}\) to \(S\). Lemma~\ref{lem:scale} supplies infinitely
many remaining primes. Fix one of them and choose a valuation
vector \(w_N=(\mu,\nu)\) from Lemma~\ref{lem:parity}.
Lemma~\ref{lem:independent} realizes that vector on the exact
scale~\eqref{eq:scale}, with the required global parameter
conditions and a unit residual. Lemma~\ref{lem:parity} then
ensures that \(\Phi\) is defined and nonzero and has even
valuation. Letting \(N\) grow makes both depths arbitrarily
large.
\end{proof}

\section{Exact examples and scope}

Take
\[
                    p=76213,\qquad B=3,\qquad n=1,\qquad T=57/4.
\]
The last value is the trace of the abscissa of the point \(P\)
in~\cite{Conditional13}. The residue \(\xi=5598\) is a simple
root of \(a(X)-1\) modulo \(p\):
\[
             \beta(\xi)\equiv28125,\qquad
             a'(\xi)\equiv22560\pmod p.
\]
For the product \(\Phi(Z_0,Z_1)=Z_0Z_1\), set
\[
             X=\frac{6603791433}{2},\qquad
             \lambda=3821558707^2+3^2.
\]
Then the valuations of \(R(\lambda)\), \(R(a(X)\lambda)\),
their product, and \(\beta(X)-1\), in that order, are
\[
                              (1,1,2,0).
\]
For the harmonic combination
\(\Phi(Z_0,Z_1)=Z_0Z_1/(Z_0+Z_1)\), set
\[
             X=\frac{24884609587871}{2},\qquad
             \lambda=304951751852576^2+3^2.
\]
The corresponding valuations are
\[
                              (1,2,2,0).
\]
The supplied exact-arithmetic checker verifies primality,
the root and derivative residues, the norm formulas, and all
four valuations in both examples. These finite checks
illustrate the theorem; the universal assertion is proved
in the preceding sections. The choice \(B=3\) in the examples
is a coefficient fixture, not a claim that this shift certifies
the thirteen-variable definition. The theorem applies separately
to every fixed odd \(B\ge3\).

The obstruction is restricted to replacing the two contact
values by one fixed nonzero rational function of those values
and demanding that its odd valuation detect their union
uniformly. It includes the product and the harmonic combination,
and it persists under the scale and positivity conditions
above. It does not assert that the constructed nonintegral
parameters satisfy the elliptic tests or the pole equations of
the full Diophantine source. It does not exclude a different
height estimate for a rational combination, a construction using
additional arithmetic conditions, functions with further
independent arguments, or a definition with eleven or fewer
variables.

\begin{thebibliography}{9}
\bibitem{Conditional13}
\emph{Thirteen Variables from a Weak \(abc\) Inequality},
9 November 2026, project paper,
\href{https://github.com/JoseBrox/diophantine-z-q_improvements/blob/main/papers/conditional_few_variables/conditional13_20261109_v16.tex}
{\texttt{conditional13\_20261109\_v16.tex}}.
\end{thebibliography}
\end{document}
